\chapter{Integration over Sections}


\newcommand{\starG}{\star}     % Hodge star in the Euclidean plane
\newcommand{\qvec}{\mathbf w}  % = \kvec\times\rvec

\newcommand{\astar}{\!*\,}             % 2-D Hodge star (acts as 90° rotation)


\newcommand{\dd}{\mathrm d}
\newcommand{\ii}{\iota}
\newcommand{\x}{\mathbf x}
\newcommand{\kvec}{\mathbf i}      % = \FrameBasis
\newcommand{\starE}{\star} % Euclidean Hodge star in 2-D

% \newcommand{\interior}[2]{#1\, \mathbin{\lrcorner}\, #2}
\newcommand{\interior}[2]{{\iota}_{#1} {#2}}

\newcommand{\IntegrationNotationTable}{

\begin{table}[H]
\centering 
\begin{tabular}{rlll}
$\WarpSV$ & $\in \Omega^{0}(\PlaneManifold)$ & A scalar field over \(\PlaneManifold\) \\
$\dd\WarpSV$
  & $\in \Omega^1(\PlaneManifold)$
  & the exterior derivative of the scalar field $\WarpSV$ \\
$\interior{\bm{\nu}}{\dd \WarpSV}$ 
  & $\in\Omega^{0}(\PlaneManifold)$ 
  & \\
$\mu$
  & $\in \Omega^2(\PlaneManifold)$
  & A volume form \\
$\omega \triangleq \interior{\FrameBasis\times \PlanePosition}{\dd A}$
  & $\in \Omega^1(\PlaneManifold)$
  & \\
$\dd \omega$
  & $\in \Omega^2(\PlaneManifold)$
  & \\
$\varrho$ 
  & $\in \Omega^1(\PlaneManifold)$ 
  & \\
\end{tabular}
\end{table}

where \(\Omega^k(\PlaneManifold)\) is the space of \(k\)-forms on \(\PlaneManifold\).

}
\subsection{Preliminaries}

We recall the following results on a Riemanian manifold \(\PlaneManifold\) with oriented volume form \(\mu\) \citep{marsden1994mathematical,abraham1988manifolds}.

\begin{ambox}[Calculus on Manifolds]{box:recall}
Definitions:
\begin{enumerate}
\item \textbf{Interior product} 
$\interior{\bm{X}}{}\colon \Omega^p(\PlaneManifold) \rightarrow \Omega^{p-1}(\PlaneManifold)$ with a vector field \(\bm{X}\)  is 
the {unique graded derivation of degree \(-1\)} satisfying  
\begin{equation}
\interior{\bm{X}}{f} = 0,
\qquad
\interior{\bm{X}}{\omega} = \omega(\bm{X})
\qquad
f\in\Omega^{0}(\PlaneManifold),
\;\omega\in\Omega^{1}(\PlaneManifold)
\label{eq:def-interior-base}
\end{equation}
and the graded-Leibniz rule 
\[
\interior{\bm {X}}{(\alpha\wedge \eta)}
      = (\interior{\bm X}{\alpha})\wedge \eta
      + (-1)^{k}\,\alpha\wedge\interior{\bm X}{ \eta},
\qquad
\alpha\in\Omega^{k} (\PlaneManifold)
\]

% \begin{itemize}
%     \item \(\interior{\bm{X}}{f} = 0\) for a scalar (0-form) field \(f\)
%     \item \(\interior{\bm{X}}{\omega} = \omega(\bm{X})\) for a $1$-form \(\omega\)
% \end{itemize}
\item \textbf{Exterior product}
$\displaystyle 
\wedge : \Omega^{k} (\PlaneManifold)\times\Omega^{\ell} (\PlaneManifold)\;\longrightarrow\;\Omega^{k+\ell} (\PlaneManifold)
$
is a bilinear, associative product defined by % the graded-anticommutativity
\[
\alpha\wedge \eta \;=\; (-1)^{k\ell}\, \eta\wedge\alpha,
\qquad
\alpha\in\Omega^{k} (\PlaneManifold),\; \eta\in\Omega^{\ell} (\PlaneManifold),
\]
with \(f\wedge\alpha = f\,\alpha\) for functions \(f\in\Omega^{0} (\PlaneManifold)\).
\item \textbf{Hodge operator} 
$\displaystyle 
\star \;:\;\Omega^{k} (\PlaneManifold)\longrightarrow\Omega^{n-k} (\PlaneManifold)
$
is the unique map satisfying  
\begin{equation}
\alpha\wedge\star \eta \;=\;
\langle\alpha, \eta\rangle\,\mu,
\qquad\forall\,\alpha, \eta\in\Omega^{k} (\PlaneManifold)
\label{eq:hodge-defn}
\end{equation}
Consequences include (with the musical isomorphism  ``\(\flat\)'' lowering indices):
\[
\interior{\bm{X}}{\star \omega} = (-1)^k \star (\bm{X}^{\flat} \wedge \omega)
\qquad\text{for}\qquad \omega \in \Omega^k(\PlaneManifold)
\]
and the special case when \(\alpha=1\), \(\star \alpha = \mu\)
\begin{equation}
\interior{\bm{X}}{\mu} = \star \bm{X}^{\flat}
\label{eq:hodge-zero}
\end{equation}
% \[
% \iota_X\mu=\star X^{\flat},
% \qquad
% \iota_X\star\omega
%    =(-1)^{k}\;\star\!\bigl(X^{\flat}\wedge\omega\bigr).
% \]


% \item \textbf{Hodge star} $\star: \bigwedge^k V \rightarrow \bigwedge^{n-k} V$
% $$
% (\alpha \wedge \star  \eta)_p=\left\langle\alpha_p,  \eta_p\right\rangle \mu_p, \quad \forall p \in M,
% $$
% for any pair of smooth $k$-forms $\alpha,  \eta \in \Omega^k (\PlaneManifold)$.
% We have the identity:
% \[
% \interior{\bm{X}}{\star \omega} = (-1)^k \star (\bm{X}^{\flat} \wedge \omega)
% \qquad\text{for}\qquad \omega \in \Omega^k(\PlaneManifold)
% \]
% and the special case when \(\alpha=1\), \(\star \alpha = \mu\)
% \[
% \interior{\bm{X}}{\mu} = \star \bm{X}^{\flat}
% \]
\item \textbf{Exterior derivative} 
$\displaystyle 
\dd : \Omega^{k} (\PlaneManifold)\longrightarrow\Omega^{k+1} (\PlaneManifold)
$
is the unique linear operator of degree \(+1\) satisfying  
\[
\begin{aligned}
\dd f\,(X)                &= X[f], && f\in\Omega^{0} (\PlaneManifold), \\[2pt]
\dd(\alpha\wedge \eta)    &= \dd\alpha\wedge \eta + (-1)^{k}\,\alpha\wedge\dd \eta,
&& \alpha\in\Omega^{k} (\PlaneManifold),\; \eta\in\Omega^{\ell} (\PlaneManifold), \\[2pt]
\dd\!\circ\dd             &= 0.
\end{aligned}
\]
Useful relations include 
\begin{equation}
    \dd f = (\nabla f)^{\flat}
    \label{eq:ident-ext-d}
\end{equation}
%
%
\item \textbf{Lie derivative} is defined on forms as the operator:
\begin{equation}
\mathcal{L}_{\bm{X}} \;=\;\interior{\bm{X}}{\dd} + \dd\interior{\bm{X}}{}
\label{eq:def-lie-form}
\end{equation}
\item \textbf{Divergence} 
\(\operatorname{div}_{\mu}\bm{X}\)
of a vector field \(\bm{X}\) with respect to a volume form \(\mu\)
\begin{equation}
(\operatorname{div}_{\mu}\bm{X})\mu \triangleq \mathcal{L}_{\bm{X}}\mu
\label{eq:def-div}
%\equiv \dd (\interior{\bm{X}}{\mu})
\end{equation}
This furnishes the identities (for a \emph{positive} scalar field \(\alpha\))
\[
(\operatorname{div}_{\mu}\bm{X})\mu = \dd (\interior{\bm{X}}{\mu})
\qquad\text{and}\qquad
\alpha \operatorname{div}_{\alpha \mu} \bm{X} = \operatorname{div}_{\mu}(\alpha \bm{X})
\]
\end{enumerate}
% Identities:
% \begin{enumerate}
%     \item $\nabla_\nu \varphi=\nu(\varphi) = \dd \varphi(\nu)=\iota_\nu \dd \varphi$
% \end{enumerate}

Theorems:
\begin{enumerate}
\item \textbf{Stokes-Cartan Theorem}
and the \textbf{Divergence Theorem} as a special case:
\begin{equation}
\int_{\PlaneManifold} \dd \alpha = \int_{\PlaneBoundary} \alpha
\qquad\text{and}\qquad
\int_{\PlaneManifold} (\operatorname{div}_{\mu} \bm{X}) \, \mu = \int_{\PlaneBoundary} \interior{\bm{X}}{\mu}
= \int_{\PlaneBoundary}\langle \bm{X}, \bm{\nu}\rangle \eta_{\mu}
\label{eq:gauss-div-thm}
\end{equation}
where $\bm{\nu}$ is the unit vector normal to $\PlaneBoundary$. 
Here we have the identities:
$$
\interior{\bm{X}}{\mu} = \langle \bm{X},\, \bm{\nu}\rangle \underbrace{\interior{\bm{\nu}}{\mu}}_{:= \eta_{\mu}}
$$
\end{enumerate}
\end{ambox}


% \iffalse
Define the 2-forms:
\[
\begin{aligned}
    \varrho &:= \interior{\nabla\WarpFA}{\volG}
            = G\,\starE \dd\WarpFA \;\in\Omega^{1}\PlaneManifold
    &\qquad
    \omega &:= \ii_{\qvec}\dd A
            = \qvec^{\flat},
\end{aligned}
\]
where ``\(\flat\)'' lowers indices \citep{abraham1988manifolds}
The Saint–Venant relation
\(
   \operatorname{div}\bigl(\WeightG(\nabla\WarpSV+\qvec)\bigr)=0
\)
is
\begin{equation}
   \dd \,\Bigl(\ii_{\nabla\WarpSV+\qvec}\volG\Bigr)=0
   \quad\text{in }\PlaneManifold,
   \qquad
   \ii_{\nu}(\dd\WarpSV+\omega)=0
   \quad\text{on }\PlaneBoundary
   % \tag{2}
   \label{eq:warp-sv-bvp}
\end{equation}

\[
% \boxed{\;
   \dd\bigl(\starE(G(\dd\WarpSV+\omega))\bigr)=0
   \quad\text{in }\Omega,
   \qquad
   \ii_{\nu}(\dd\WarpSV+\omega)=0
   \quad\text{on }\partial\Omega.
% \;}
% \tag{2}
\]
% \fi
%
%
%

\iffalse
\subsubsection{Poisson Scaling}
Fix a normalization $\int_{\mathscr D} u_c=0$ for each $c$, and write the solution family as

$$
u_c = u_0 + c\,w,
$$

where

$$
-\Delta u_0=f,\ \ \partial_n u_0=0,\qquad
-\Delta w=0,\ \ \partial_n w=v\!\cdot n,
$$

all on the same domain $\mathscr D$.

\begin{enumerate}

\item $\displaystyle \mathtt{cnv} \approx \int_{\mathscr D}\nabla u_c\,dx$

Use Green’s first identity with the coordinate test functions $\phi=y$ and $\phi=z$:

$$
\int_{\mathscr D}\nabla u_c\cdot \nabla \phi\,dx
= \int_{\partial\mathscr D}\phi\,\partial_n u_c\,ds
  - \int_{\mathscr D}\phi\,\Delta u_c\,dx.
$$

Since $\nabla \phi=e_y$ or $e_z$, this gives 
% componentwise

% $$
% \int_{\mathscr D}\partial_y u_c\,dx
% = c\int_{\partial\mathscr D} y\,(v\!\cdot n)\,ds\;+\;\int_{\mathscr D} y\,f\,dx,
% $$

% $$
% \int_{\mathscr D}\partial_z u_c\,dx
% = c\int_{\partial\mathscr D} z\,(v\!\cdot n)\,ds\;+\;\int_{\mathscr D} z\,f\,dx.
% $$

% Vectorially,

$$
\boxed{\ \int_{\mathscr D}\nabla u_c\,dx
= c\int_{\partial\mathscr D} \bm{r}\,(\bm{v}\!\cdot \bm{n})\,ds\;+\;\int_{\mathscr D} \bm{r}\,f\,dx\ ,}
$$

Hence this integral depends \emph{affinely} on $c$ (linear slope plus a constant offset from $f$).

\item $\displaystyle \int_{\mathscr D}\nabla u_c\otimes \nabla u_c\,dx$

Expand using $u_c=u_0+c w$:

$$
\int_{\mathscr D}\nabla u_c\otimes \nabla u_c\,dx
= A \;+\; c\,S \;+\; c^2\,C,
$$

where

$$
A:=\int_{\mathscr D}\nabla u_0\otimes\nabla u_0\,dx,\quad
C:=\int_{\mathscr D}\nabla w\otimes\nabla w\,dx\ \ (\text{sym.\ \& psd}),
$$

$$
S:=\int_{\mathscr D}\!\big(\nabla u_0\otimes\nabla w+\nabla w\otimes\nabla u_0\big)\,dx\ \ (\text{sym.}).
$$
\end{enumerate}


So this matrix-valued integral is a \emph{quadratic polynomial in $c$}. 
In particular, its trace (the Dirichlet energy) satisfies

$$
\int_{\mathscr D}|\nabla u_c|^2\,dx
= \int_{\mathscr D}|\nabla u_0|^2\,dx
+ 2c\int_{\mathscr D}\nabla u_0\cdot \nabla w\,dx
+ c^2\int_{\mathscr D}|\nabla w|^2\,dx,
$$

and the mixed/curvature coefficients admit boundary representations via Green’s identity:

$$
\int_{\mathscr D}\nabla u_0\cdot \nabla w\,dx
= \int_{\partial\mathscr D} u_0\,(v\!\cdot n)\,ds,\qquad
\int_{\mathscr D}|\nabla w|^2\,dx
= \int_{\partial\mathscr D} w\,(v\!\cdot n)\,ds.
$$

\paragraph{Summary}
With mean-zero normalization and free $c$:
\begin{itemize}
\item $\displaystyle \int_{\mathscr D}\nabla u_c$ is \emph{affine} in $c$ with slope $ \int_{\partial\mathscr D} r\,(v\!\cdot n)\,ds$.
\item $\displaystyle \int_{\mathscr D}\nabla u_c\otimes\nabla u_c$ is \emph{quadratic} in $c$ with coefficients determined by $u_0$ and the harmonic response $w$ to the unit flux $v\!\cdot n$.
\end{itemize}
\fi

\subsection{Identities}


%
% \cite{armero2021modeling}
% \begin{equation*}
% \SectionIntegral{\nabla \WarpSV + \FrameBasis \times (\PlanePosition - \tilde{\bm{\SectionLocation}}) \volG} = 0 \\
% \quad\implies\quad
%     \SectionIntegral{\nabla \WarpSV \volG} = \SectionIntegral{\PlanePosition\times\FrameBasis \volG}
% \end{equation*}
% \cite{armero2021modeling} Equation~(102) and (103)
% \begin{equation*}
%     \Jv = \SectionIntegral{\nabla \WarpSV \cdot \nabla \WarpSV \, \volG} 
%     = - \SectionIntegral{ [\nabla \WarpSV, \FrameBasis, \PlanePosition ] \, \volG} 
%     > 0
% \end{equation*}

\begin{proposition}
\label{prop:div-skew}
Let $(\PlaneManifold,\mu)$ be a compact oriented manifold with (possibly non-empty) boundary $\PlaneBoundary$.  
If a vector field $\bm{X}$ is $\mu$-divergence free in the sense:
\[
\operatorname{div}_{\mu} \bm{X} = 0
\quad\text{and}\quad
\bm{X}\cdot\bm{\nu} = 0 \;\; \text{on } \PlaneBoundary,
\]
then $\bm{X}$ acts as a skew–symmetric operator on $\mathcal{F}(\PlaneManifold)$:
\[
% \int_{\PlaneManifold} \bm{X}[f]\,g\,\mu \;=\; - \int_{\PlaneManifold} f\,\bm{X}[g]\,\mu ,
\left< \bm{X}[f],\,g\right>_\mu \;=\; - \left<f,\,\bm{X}[g]\right>_{\mu} ,
\qquad\forall\,f,g\in\mathcal F(\PlaneManifold).
\]
with the operation \(\bm{X}[f] \triangleq \nabla_{\bm{X}} f\) .
\end{proposition}

\begin{proof}
This is an extension of \cite[Corollary 7.2.11]{abraham1988manifolds}. 
Because $\operatorname{div}_{\mu}\bm{X}=0$, for any $h\in\mathcal F(\PlaneManifold)$ we have
\[
\mathcal{L}_{\bm{X}}(h\,\mu)=\bm{X}[h]\,\mu.
\]
Taking $h=fg$ gives
\[
\bm{X}[fg]\,\mu
= \mathcal{L}_{\bm{X}}(fg\,\mu).
\]
Integrating over $\PlaneManifold$ and applying the divergence theorem yields
\[
\int_{\PlaneManifold} \bm{X}[fg]\,\mu
   \;=\;
\int_{\PlaneBoundary} (fg)\,(\bm{X}\cdot\bm{\nu})\,\eta
\]
Since $\bm{X} \cdot \bm{\nu}=0$ on $\PlaneBoundary$, the integrals vanish and
\[
0
=\int_{\PlaneManifold} \bm{X}[fg]\,\mu
  =\int_{\PlaneManifold} \bigl(\bm{X}[f]\,g + f\,\bm{X}[g]\bigr)\mu.
\]
\end{proof}

\begin{proposition}
\label{prop:7212}
\citep[Corollary 7.2.12]{abraham1988manifolds}
If \(\PlaneManifold\) is compact without boundary, $\bm{X} \in \mathfrak{X}(\PlaneManifold), \alpha \in \Omega^{\mathrm{k}}(\PlaneManifold)$ and $\beta \in \Omega^{\mathrm{n}-\mathrm{k}}(\PlaneManifold)$, then

$$
\int_{\PlaneManifold} \mathcal{L}_{\bm{X}} \alpha \wedge \beta=-\int_{\PlaneManifold} \alpha \wedge \mathcal{L}_{\bm{X}} \beta
$$

\end{proposition}
\begin{proof}
Since $\alpha \wedge \beta \in \Omega^n(M)$, the formula follows by integrating both sides of the relation 
\[
\dd \interior{\bm{X}}{(\alpha \wedge \beta)}=\mathcal{L}_X(\alpha \wedge \beta)=\mathcal{L}_X \alpha \wedge \beta+\alpha \wedge \mathcal{L}_X \beta
\]
and using \eqref{eq:gauss-div-thm} (Stoke's Theorem).
\end{proof}

\iffalse
\subsubsection*{Riesz I (81b)}
\fi

% $$
% \int_{\PlaneManifold} \nabla \WarpFA \cdot \nabla \WarpSV \; \volG
% =
% \int_{\PlaneManifold} \iota_{\nabla \WarpFA}( \dd \WarpSV) \wedge \mu
% $$

\iffalse
\subsection{Classical}

\begin{enumerate}
\item For any vector field \(\bm w\) on \(\Section\), Green’s identity applied componentwise yields
\begin{equation}
\int_{\Section}\nabla\bm w\,(\nabla\bm w)^{\!\top}\,\dA
= -\int_{\Section}\bm w\otimes \Delta\bm w\,\mathrm dA
\;+\;\int_{\SectionBoundary}\bm w\otimes \partial_{\bm\nu}\bm w\,\mathrm ds.
\label{eq:thm-green-vec}
\end{equation}
\item Also, by the divergence theorem,
\begin{equation}
\int_{\Section}\nabla\bm w\,\dA
=\int_{\SectionBoundary}\bm\nu\otimes\bm w\,\mathrm ds.
\label{eq:div-thm-vec}
\end{equation}
\item Rotated-gradient integration by parts (Stokes/Green in 2-D). For any vector field \(\bm a\) and scalar \(g\),
\begin{equation}
\int_{\Section}\bm a\otimes\FrameBasis\times\nabla g\,\dA
=\oint_{\partial\Section} g\,\bm a\otimes(\bm\nu\times\FrameBasis)\,\mathrm ds
-\int_{\Section} g\,(\nabla\bm a)\,J^{\top}\,\dA.
\label{eq:thm-stokes-green}
\end{equation}
\end{enumerate}
\fi 

\subsection{Torsion}
\begin{proposition}
\label{prop:armero-81}
The solution \(\WarpSV\) to the boundary value problem \eqref{eq:warp-sv-bvp} is the Riesz vector:
\begin{equation}
    (\WarpSV,\, \phi)_G =  \frac{1}{2} \left\{\PlanePosition \cdot \PlanePosition,\, \phi\right\}_G %= - [\nabla \phi,\, \FrameBasis, \PlanePosition]_G =
    \qquad \forall \phi \in \mathcal{F}(\PlaneManifold)
\end{equation}
where \(\{\,\cdot,\,\cdot\,\}_G\) is the integrated Poisson bracket \citep{arnold1978mathematical} given by 
\[
\left\{\psi,\, \phi\right\}_G = \int_{\PlaneManifold} \nabla \psi \cdot \mathbb{J} \nabla \phi \, \mu_G
\]
and the symplectic operator is \(\mathbb{J} = \FrameBasis^{\times}\).
This gives the special cases 
%
\begin{equation}
\int [ \nabla \WarpFA,\; \FrameBasis,\; \PlanePosition - \tilde{\bm{\SectionLocation}} ] \, \volG
= \int_{\PlaneManifold} \WarpSV^2 \, \volE % = \Jw
\label{eq:FA81b}
\end{equation}
and
\begin{equation}
% \Jv := 
\int_{\PlaneManifold} \nabla \WarpSV \cdot \nabla \WarpSV \; \volG 
= - \int_{\PlaneManifold} [\nabla \WarpSV, \FrameBasis, \PlanePosition] \; \volG > 0
\label{eq:FA103}
\end{equation}
\end{proposition}
\begin{proof}
By \cref{prop:div-skew}, taking \(\bm{X} = \nabla \WarpSV + \FrameBasis \times \PlanePosition\), \(f = \phi\) and \(g = 1\).
\end{proof}

\iffalse
\subsubsection*{Riesz II (88/81a)}
\fi 

\begin{proposition}
\label{prop:FA88}
\cite{armero2021modeling} Equation~(88)
\begin{equation}
(\WarpFA,\, \phi)_G = - \left<\WarpSV,\, \phi\right>_E
\qquad \forall \phi \in \mathcal{F}(\PlaneManifold)
\label{eq:FA88}
\end{equation}
\end{proposition}
% \begin{equation}
%    \int_{\PlaneManifold}
%       (\nabla\WarpFA\cdot\nabla\WarpSV)\,G\,\dd A
%    =\int_{\PlaneManifold} \WarpSV^{2}\,E\,\dd A.
% \label{eq:FA-SV2}
% \end{equation}
\begin{proof}
Multiply the PDE for $\WarpFA$ by $\phi$, and integrate over $\PlaneManifold$

$$
\int_{\PlaneManifold} \phi \, (\operatorname{div}_{G}\nabla \WarpFA)  \volG 
= \int_{\PlaneManifold} \phi \, \WarpSV \volE
% = \int_{\PlaneManifold} \WarpSV \, \frac{n_E}{n_G} \WarpSV  \mu
$$
We thus aim to show that the left hand side is equivalent to the (negative) $G$-energy inner product:
$$
\int_{\PlaneManifold} \phi\, (\operatorname{div}_{G} \nabla \WarpFA) \volG
=- \int_{\PlaneManifold} \nabla \WarpFA \cdot \nabla \phi \, \volG 
$$
This is equivalent to showing:
\begin{equation}
\int_{\PlaneManifold} \phi\, (\operatorname{div}_{G} \nabla \WarpFA) \volG
+ \int_{\PlaneManifold} \dd \phi \wedge \interior{\nabla \WarpFA}{\volG} = 0
\tag{$\dagger$}
\end{equation}
where we have used 
\[
\begin{array}{rlr}
\dd \phi \wedge \interior{\nabla \WarpFA}{\volG} 
&= \dd \phi \wedge \star (\nabla \WarpFA)^{\flat}  & \qquad \text{\eqref{eq:hodge-zero}}\\
&= \dd \phi \wedge \star     \dd \WarpFA &  \qquad \text{\eqref{eq:ident-ext-d}} \\
&= \left<\dd \WarpFA,\, \dd \phi \right> \volG & \qquad \text{\eqref{eq:hodge-defn}} \\
\end{array}
\]
Thus (\(\dagger\)) is exactly an integration by parts of the product rule for \eqref{eq:def-lie-form}
\[
\mathcal{L}_{(\WarpSV \nabla \WarpFA)} \volG = \WarpSV \mathcal{L}_{\nabla \WarpFA}\volG + \dd \WarpSV \wedge \interior{\nabla \WarpFA}{\volG}
\]
where by the divergence theorem \eqref{eq:gauss-div-thm}
\[
\int_{\PlaneManifold} \mathcal{L}_{(\WarpSV \nabla \WarpFA)} \volG =
\int_{\PlaneBoundary} \WarpSV  \nabla \WarpFA \cdot \bm{\nu} \; \eta_G
\]
which vanishes due to the boundary condition \(\nabla_{\bm{\nu}} \WarpFA = 0\). 
Therefore:

$$
-\int_{\PlaneManifold} \nabla \WarpFA \cdot \nabla \WarpSV \, \volG
= \int_{\PlaneManifold} \WarpSV^2 \volE
$$
and also
$$
\int_{\PlaneManifold} \interior{\nabla\WarpFA}{\dd \WarpSV}\,\volG
=
- \int_{\PlaneManifold} \WarpSV^2 \; \volE
$$

\end{proof}
\begin{proof}[Alternative Proof]
Multipling the PDE for $\WarpFA$ by $\phi$, and integrating over $\PlaneManifold$ yields

$$
\int_{\PlaneManifold} \phi \, (\operatorname{div}_{G}\nabla \WarpFA)  \volG 
= \left<\phi,\, \WarpSV\right>_E
% = \int_{\PlaneManifold} \WarpSV \, \frac{n_E}{n_G} \WarpSV  \mu
$$
We thus aim to show that the left hand side is equivalent to the (negative) $G$-energy inner product:
$$
\int_{\PlaneManifold} \phi\, (\operatorname{div}_{G} \nabla \WarpFA) \volG
=- \int_{\PlaneManifold} \nabla \WarpFA \cdot \nabla \phi \, \volG 
$$
Applying \cref{prop:7212} with \(\alpha = \phi \in \Omega^0(\PlaneManifold)\), \(\beta = \volG \in \Omega^2(\PlaneManifold)\), and \(\bm{X} = \nabla \WarpFA\) furnishes
\[
\int_{\PlaneManifold} \phi \wedge \underbrace{(
\operatorname{div}_G \nabla \WarpFA
)\volG}_{\mathcal{L}_{\bm{X}}\volG}  = -\int_{\PlaneManifold} (\mathcal{L}_{\nabla \WarpFA} \phi ) \wedge \volG 
\]
Recall that the definitions for the Lie derivative \(\mathcal{L}_{\bm{X}} \phi\), exterior derivative \(\dd \phi\), and directional derivative \(\bm{X}[\phi] = \nabla_{\bm{X}} \phi\) on a 0-form \(\phi \in \Omega^0(\PlaneManifold)\) all coincide, so the left side of the equality is rewritten as:
\[
\int_{\PlaneManifold} (\underbrace{\nabla_{\nabla \WarpFA} \phi}_{\equiv \nabla \WarpFA \cdot \nabla \phi}) \wedge \volG 
= - \int_{\PlaneManifold} \phi \wedge (\operatorname{div}_G \nabla \WarpFA)\volG 
\]
Recall also that the wedge product on scalar fields degenerates into scalar multiplication.
\end{proof}

\iffalse
\begin{remark}
Using integration by parts via Generalized Stokes to move one derivative off of $\dd \WarpSV$, using the \emph{codifferential} $\delta$, the adjoint of $\dd$, defined by:

$$
\int_{\PlaneManifold} \langle \dd\WarpFA, \dd\WarpSV \rangle \, \volG = \int_{\PlaneManifold} \WarpFA \cdot \delta \dd\WarpSV \, \volG + \int_{\PlaneBoundary} \WarpFA \cdot \interior{\bm{\nu}}{} \dd\WarpSV \, \mu_{\partial}
$$

% Let $\phi = \WarpFA$, $\phi = \WarpSV$ and $\delta \dd \WarpSV = \operatorname{div}_\mu(\nabla \WarpSV)$
% so
% $$
% \int_{\PlaneManifold} \langle \dd\WarpFA, \dd\WarpSV \rangle \, \mu
% = -\int_{\PlaneManifold} \WarpFA \cdot \delta d\WarpSV \, \mu
% + \int_{\PlaneBoundary} \WarpFA \cdot \iota_\nu \dd\WarpSV \, \sigma_G
% $$

Now, since $\delta \dd\WarpSV = \operatorname{div}_{G} (\nabla \WarpSV)$, and assuming regularity, we have:

$$
\int_{\PlaneManifold} \nabla \WarpFA \cdot \nabla \WarpSV \; \volG
=
- \int_{\PlaneManifold} \WarpFA  \operatorname{div}_\mu (\nabla \WarpSV) \; \volG
+ \int_{\PlaneBoundary} \WarpFA  \nabla \WarpSV \cdot \bm{\nu} \; \sigma_G
$$
or

$$
\int_{\PlaneManifold} \nabla \WarpFA \cdot \nabla \WarpSV \; \volG
+ \WarpFA  \operatorname{div}_\mu (\nabla \WarpSV) \; \volG 
=
\int_{\PlaneBoundary} \WarpFA  \nabla \WarpSV \cdot \bm{\nu} \; \sigma_G
$$
where $\sigma_G = n_G \, \dd s$:
\end{remark}
\fi

\iffalse
\begin{proposition}[Relocation]
    Consider the BVP:
    \begin{equation}
    \left\{
    \begin{array}{rll}
    \Delta \varphi_{\SectionLocation} &=0 & \mathrm{in }\ \PlaneManifold, \\
    \nabla_{\mathbf{n}} \varphi_{\SectionLocation}  &=-\FrameBasis \times (\bm{r}-\bm{\SectionLocation}) \cdot \mathbf{n} & \mathrm{ on }\ \partial \PlaneManifold  \\
    \SectionIntegral{\varphi_{\SectionLocation}} &= 0
    \end{array}
    \right.
    \label{eq:torsion-neumann}
    \end{equation}

    If \(\varphi_a\) is the solution to \eqref{eq:torsion-neumann} with \(\bm{\pi} = \bm{\SectionLocation}_a\), then 
    \[
    \varphi_{b} \triangleq \varphi_a + [\FrameBasis, \bm{\SectionLocation}_{ab}, \bm{\SectionCoordinate}] - [\FrameBasis, \bm{\SectionLocation}_{ab},  \bar{\bm{\SectionLocation}}] % TODO: A?
    \]
    is the solution for \(\bm{\SectionLocation}_{b} \triangleq \bm{\SectionLocation}_a + \bm{\SectionLocation}_{ab}\)
\end{proposition}
\begin{proof}
Our objective is to show that the function \(\varphi_b\) (i) is harmonic, (ii) satisfies the the boundary condition \eqref{eq:torsion-neumann}$_2$, and (iii) satisfies the normalizing condition.
\begin{equation*}
\nabla_{\mathbf{n}}\varphi_b = \nabla_{\mathbf{n}}\varphi_a + \nabla_{\mathbf{n}}\left[\FrameBasis,\, \bm{\SectionLocation}_{ab},\, \SectionCoordinate \right]
\end{equation*}

\end{proof}
\fi